\begin{table}[H]
\centering
\small
\caption{Distribution of pairs and regression weight by plate appearances per pair}
\label{tab:pair_pa_distribution}
\begin{tabular}{lrrrr}
\toprule
PAs per pair & Pairs & \% of pairs & PAs (weight) & \% of weight \\
\midrule
\multicolumn{5}{l}{\textit{Panel A: 2025}} \\
1--30 & 8,334 & 84.7 & 63,266 & 35.5 \\
31--50 & 679 & 6.9 & 26,468 & 14.9 \\
51--75 & 352 & 3.6 & 21,806 & 12.2 \\
76--100 & 181 & 1.8 & 15,683 & 8.8 \\
101+ & 288 & 2.9 & 50,809 & 28.5 \\
\addlinespace
Total & 9,834 & 100.0 & 178,032 & 100.0 \\
\midrule
\multicolumn{5}{l}{\textit{Panel B: 2026}} \\
1--30 & 8,481 & 84.6 & 65,497 & 36.6 \\
31--50 & 715 & 7.1 & 27,700 & 15.5 \\
51--75 & 398 & 4.0 & 24,122 & 13.5 \\
76--100 & 159 & 1.6 & 13,819 & 7.7 \\
101+ & 275 & 2.7 & 47,853 & 26.7 \\
\addlinespace
Total & 10,028 & 100.0 & 178,991 & 100.0 \\
\bottomrule
\end{tabular}
\par\smallskip\parbox{0.9\textwidth}{\footnotesize \textit{Notes:} Each observation is a (focal hitter, conditional hitter) pair within a season. PAs per pair is the number of plate appearances in which the focal hitter batted directly behind the conditional hitter. The regressions weight each pair by this number, so the weight share equals the share of total pair PAs. Hitters with fewer than 50 PA in a season are excluded, and no minimum is imposed on PAs per pair.}
\end{table}
